\chapter{Expanders and Ramanujan Graphs}\label{ch:expanders}

\begin{center}\itshape
A sparse network can have no cheap cut. Its eigenvalues tell us when,\\
and a zeta function records the same constraint through its poles.
\end{center}

Give every vertex only a fixed number of neighbours. Can the resulting network still connect large pieces so thoroughly that no small collection of edges separates them? The preceding chapter answered this through a spectral gap. Now we ask how to build such graphs, what a walk gains from the gap, and why the number $2\sqrt{d-1}$ keeps appearing.

Two 1973 papers approached sparse networks differently. Pinsker proved the existence of concentrators using randomness; Margulis gave an explicit construction using group actions. An existence argument and an explicit construction solve different problems: the former says a suitable graph is available, while the latter specifies one. Later constructions by Lubotzky, Phillips and Sarnak, and independently Margulis, reached the tree's spectral bound using number theory. Their name, \emph{Ramanujan graphs}, recalls the modular-form eigenvalue estimate behind the construction.

The tree gives both a counting problem and a limiting distribution. Its return-walk counts are moments of the Kesten--McKay law. Large graphs whose fixed-radius neighbourhoods are mostly trees have the same limiting law. Closed walks with immediate reversals forbidden give a different count, assembled by Ihara into a zeta function. Bass's determinant formula then turns its poles back into adjacency eigenvalues. The Ramanujan bound becomes a Riemann hypothesis for that zeta function, with a corresponding estimate for primitive cycles.

There is a recent answer to how often the bound holds. Huang, McKenzie and Yau's December 2024 preprint proves that, for each fixed $d\ge3$, a uniformly random simple $d$-regular graph is Ramanujan with limiting probability about $69\%$ as the admissible vertex count tends to infinity. The result concerns both edges of the spectrum, not only $\lambda_2$. It settled a long-running question discussed by Alon and Sarnak; Sloman's 2025 account tells the story. We quote this result, rather than assuming its random-matrix proof.

\begin{reminder}[spectral bounds already available]
For adjacency eigenvalues $d=\lambda_1\ge\cdots\ge\lambda_n$, put $\lambda=\max(|\lambda_2|,|\lambda_n|)$. Cheeger's inequality gives $h(G)\ge(d-\lambda_2)/2$ (Theorem~\ref{spectral:thm:cheeger}), and the mixing lemma bounds $|e(S,T)-d|S||T|/n|$ by $\lambda\sqrt{|S||T|}$ (Lemma~\ref{spectral:lem:mixing}). Alon--Boppana gives $\lambda_2\ge2\sqrt{d-1}-o(1)$ for fixed degree and growing size (Theorem~\ref{spectral:thm:alonboppana}). Connectedness makes $d$ simple, and bipartiteness is equivalent to the occurrence of $-d$ (Exercise~\ref{linear:ex:bipartite}; also Exercise~\ref{ex:bipartiteeig} below).
\end{reminder}

\section{Spectral gap and mixing}

Here $G$ is finite, connected and $d$-regular with $d>0$. The usual walk has transition matrix $P=A/d$. Set $\alpha=\lambda/d$. On $\mathbf1^\perp$, the norm of $P$ is $\alpha$. Let $J=\mathbf1\mathbf1^{\mathsf T}/n$ denote projection onto constants; within this chapter $J$ includes the factor $1/n$. Subtracting the constant component gives
\begin{equation}\label{eq:pje}
P=J+E,\qquad EJ=JE=0,\qquad\|E\|=\alpha.
\end{equation}
For probability vectors $p,u$, their total variation distance is $\|p-u\|_{\mathrm{TV}}=\tfrac12\sum_x|p(x)-u(x)|$; it measures the largest difference in probabilities assigned to an event.

\begin{proposition}[mixing]\label{prop:mixing}
Starting from any probability law $p_0$, the walk's law $p_t$ satisfies
\[
\|p_t-u\|_{\mathrm{TV}}\le\frac{\sqrt n}{2}\alpha^t,
\qquad u=\mathbf1/n.
\]
If $\alpha<1$, a non-negative integer number of steps satisfying
$t\ge\log(\sqrt n/(2\varepsilon))/\log(1/\alpha)$
suffices for distance at most $\varepsilon$.
\end{proposition}
\begin{proof}
Exercise~\ref{spectral:ex:walkmixing} gives $\|p_t-u\|_2\le\alpha^t\|p_0-u\|_2\le\alpha^t$ for every initial law. Cauchy--Schwarz gives $\|p_t-u\|_1\le\sqrt n\|p_t-u\|_2$, proving the bound. Here $\alpha>0$, since $A$ has trace zero and a positive eigenvalue $d$. Taking logarithms proves the stated step count when $\alpha<1$.
\end{proof}

On a family with $\alpha\le1-\varepsilon_0$, a fixed total variation error requires $O(\log n)$ steps. This order cannot be improved for point starts and fixed degree: after $t$ steps the possible locations number at most $d^t$, so the complement of the support still has uniform mass at least $1-d^t/n$. For comparison, the lazy cycle mixes on the order of $n^2$ steps, and the lazy $k$-cube on the order of $k\log k$. Laziness is essential to these comparisons when the underlying graph is bipartite.

\section{Saving random bits}

Suppose an algorithm uses a $k$-bit seed and fails only for seeds in a set $B$ containing at most half the $2^k$ possibilities. For a one-sided error, independent repetition can fail only when every chosen seed is bad. Expander walks replace independent seeds by correlated ones while keeping that event unlikely. This use of walks goes back to Ajtai, Koml\'os and Szemer\'edi.

Place the seeds on a $d$-regular graph, choose a uniform start and take $t$ steps. If $d$ is a power of two, the walk uses exactly $k+t\log_2d$ random bits: $k$ for the first vertex, then $\log_2d$ for each uniform neighbour. For a general degree, exact uniform choices need a different bit-accounting convention; Exercise~\ref{ex:amplify} addresses this point.

\begin{proposition}[hitting property]\label{prop:hitting}
Let $B\subseteq V$ have size $\beta n$, and start the walk uniformly. For every non-negative integer $t$,
\[
\Pr(X_0,X_1,\ldots,X_t\in B)\le\beta(\beta+\alpha)^t.
\]
\end{proposition}
\begin{proof}
Let $\Pi$ multiply coordinates by $\mathbf1_B$. The sum of the probabilities of all paths staying inside $B$ is
\[
\frac1n\mathbf1_B^{\mathsf T}(\Pi P\Pi)^t\mathbf1_B.
\]
This formula also holds at $t=0$. From~\eqref{eq:pje}, $\Pi P\Pi=\Pi J\Pi+\Pi E\Pi$. The first term is the rank-one matrix $\mathbf1_B\mathbf1_B^{\mathsf T}/n$, of norm $\beta$; the second has norm at most $\alpha$. Bounding the quadratic expression by the operator norm gives
$n^{-1}\|\mathbf1_B\|^2(\beta+\alpha)^t=\beta(\beta+\alpha)^t$.
\end{proof}

When $\beta\le1/2$ and $\alpha\le1/4$, this gives error at most $(3/4)^t$. The number of random bits grows as $k+O(t)$ for a fixed degree with uniform sampling implemented correctly. Two-sided errors require control of how many bad seeds are visited, rather than whether every seed is bad. A Chernoff estimate supplies that control.

\begin{theorem}[expander Chernoff bound]\label{thm:expchernoff}
Let $f:V\to[0,1]$ have uniform mean $\mu$, and let the walk start uniformly. Set $S_t=\sum_{i=1}^tf(X_i)$. If $\alpha<1$, then for every $0<\varepsilon\le1$,
\[
\Pr(S_t\ge(\mu+\varepsilon)t)
\le\exp\left(-\frac{(1-\alpha)\varepsilon^2t}{20}\right).
\]
The same bound holds for $\Pr(S_t\le(\mu-\varepsilon)t)$.
\end{theorem}

The exponent has the same dependence on $t$ and $\varepsilon$ as Hoeffding's independent-sample bound (Volume~4, Theorem~4.2), with a loss governed by the gap $1-\alpha$. Gillman's work established concentration for expander walks; the following argument proves the particular constant stated here.

\begin{proof}
Choose $0<r\le(1-\alpha)/4$ and let $D$ be diagonal with entries $e^{rf(x)}$. Define the non-negative row vector $w_k=u^{\mathsf T}(PD)^k$. Summing path weights shows that $w_k(x)=\mathbb E[e^{rS_k}\mathbf1_{X_k=x}]$. Its mass $a_k=\sum_xw_k(x)$ is therefore the exponential moment. Split
\[
w_k=a_ku^{\mathsf T}+w_k^\perp,\qquad
\sum_xw_k^\perp(x)=0,\qquad b_k=\sqrt n\|w_k^\perp\|_2.
\]
Put $s=e^r-1$ and $m=n^{-1}\sum_xe^{rf(x)}$.

\emph{One step.} Multiplication by $P$ preserves the constant part and sends the remainder to $z=w_k^\perp P$, with zero mass and $\|z\|_2\le\alpha\|w_k^\perp\|_2$. Now multiply by $D$. The constant part has mass $ma_k$; after removing that mass its norm is at most $a_ks/\sqrt n$, because both $e^{rf(x)}$ and $m$ lie in $[1,e^r]$. The other part has mass
$\sum_xz(x)(e^{rf(x)}-1)$,
whose absolute value is at most $\sqrt n\,s\|z\|_2$. Orthogonal projection onto zero-mass rows does not increase norm, so its remaining part has norm at most $\|zD\|_2\le e^r\|z\|_2$. Hence
\[
a_{k+1}\le ma_k+s\alpha b_k,\qquad
b_{k+1}\le sa_k+e^r\alpha b_k,
\qquad a_0=1,\ b_0=0.
\]

\emph{An invariant bound.} For $r\le1$, the exponential series gives $s\le r+r^2$. Our choice of $r\le1/4$ improves this to $s\le5r/4$. It follows that $e^r\alpha\le\alpha+5\alpha(1-\alpha)/16\le(1+\alpha)/2$. Since $m\ge1$, we obtain $m-e^r\alpha\ge(1-\alpha)/2$.
Put $\kappa=2s/(1-\alpha)$ and $\Lambda=m+s\alpha\kappa$. Induction gives $a_k\le\Lambda^k$ and $b_k\le\kappa\Lambda^k$. The first recurrence preserves this pair of bounds directly; the second does so because
$\kappa(\Lambda-e^r\alpha)\ge\kappa(1-\alpha)/2=s$.

\emph{Estimate the growth.} Convexity on $[0,1]$ gives $e^{rx}\le1+xs$, so
$m\le1+\mu s\le e^{\mu r+r^2}$.
Also $s\alpha\kappa=2s^2\alpha/(1-\alpha)\le4r^2/(1-\alpha)$. As $e^{\mu r+r^2}\ge1$,
\[
\Lambda\le e^{\mu r+r^2}\left(1+\frac{4r^2}{1-\alpha}\right)
\le\exp\left(\mu r+\frac{5r^2}{1-\alpha}\right).
\]
Markov's inequality now gives
\[
\Pr(S_t\ge(\mu+\varepsilon)t)
\le\exp\left(-r\varepsilon t+\frac{5r^2t}{1-\alpha}\right).
\]
Take $r=\varepsilon(1-\alpha)/10$, within the allowed range. The exponent becomes $-(1-\alpha)\varepsilon^2t/20$. Applying the upper-tail result to $1-f$ proves the lower-tail assertion.
\end{proof}

\begin{think}
The proof follows the mass and the zero-mass remainder separately. What information would disappear if we used only $\|PD\|\le e^r$? That bound forgets both the mean $\mu$ and the contraction $\alpha$ on the remainder. The two recurrences retain them.
\end{think}

\begin{remark}
The constant $20$ leaves room for improvement. Work by Healy and by Chung, Lam, Liu and Mitzenmacher gives sharper estimates and extensions. For a nonuniform start, a common extra factor is the $L^2(u)$ norm of the density $p_0/u$; specifying the measure in that norm matters. Expander walks have been used for error amplification and pseudorandom computation. Reingold's logarithmic-space algorithm for undirected connectivity is another application of expander constructions, rather than a consequence of this concentration bound alone.
\end{remark}

\section{The spectral measure of the tree}

Fix the infinite $d$-regular tree $T_d$ and a vertex $o$. Count closed walks of length $k$ starting at $o$, calling the result $a_k$. These counts will be the moments of a probability measure. No infinite-dimensional spectral theorem is needed to identify it.

A walk is \emph{non-backtracking} when it never immediately reverses its preceding edge. Write $A_k(x,y)$ for the number of such walks of length $k$ from $x$ to $y$. All matrix products below are locally finite, so each entry involves only finitely many neighbours even on the infinite tree.

\begin{lemma}[non-backtracking polynomials]\label{lem:nbpoly}
On a finite or infinite $d$-regular graph,
\[
A_0=I,\quad A_1=A,\quad A_2=A^2-dI,\quad
A_{k+1}=AA_k-(d-1)A_{k-1}\quad(k\ge2).
\]
In particular $A_k=q_k(A)$ for monic degree-$k$ polynomials determined by $q_0=1$, $q_1=x$, $q_2=x^2-d$ and $q_{k+1}=xq_k-(d-1)q_{k-1}$ for $k\ge2$.
\end{lemma}
\begin{proof}
Append one edge to a non-backtracking length-$k$ walk. A forbidden extension reverses its last edge. Removing those last two steps leaves a length-$(k-1)$ non-backtracking walk. For $k\ge2$, each such shortened walk can be extended out and back in $d-1$ ways, since the outward edge must avoid its previous edge. When $k=1$, there is no previous edge, giving $d$ choices instead. Thus $A_kA=A_{k+1}+(d-1)A_{k-1}$ for $k\ge2$ and $A^2=A_2+dI$. Inductively each $A_k$ is a polynomial in $A$, hence commutes with it, giving the stated order of factors and the monic recurrence.
\end{proof}

On a tree a non-backtracking walk cannot repeat a vertex: a first repetition would form a cycle. Therefore
\begin{equation}\label{eq:treenb}
q_k(A)(o,o)=0\qquad(k\ge1).
\end{equation}
This vanishing determines all the ordinary return counts, because each $q_k$ has leading term $x^k$.

\begin{theorem}[Kesten, 1959]\label{thm:kestenmeasure}
For $d\ge3$, the number $a_k$ of length-$k$ closed walks from a fixed vertex of $T_d$ is
$a_k=\int x^k\nu_d(\dd x)$,
where the probability measure $\nu_d$ has density
\[
\nu_d(\dd x)=\frac{d\sqrt{4(d-1)-x^2}}{2\pi(d^2-x^2)}
\mathbf1_{|x|\le2\sqrt{d-1}}\dd x.
\]
\end{theorem}
\begin{proof}
Define the linear functional $L$ on polynomials by $L(x^k)=a_k$. Equation~\eqref{eq:treenb} gives $L(q_0)=1$ and $L(q_k)=0$ for $k\ge1$. Since the $q_k$ are monic of consecutive degrees, they form a polynomial basis. It suffices to check the same identities by integration against the proposed density.

Put $\sigma=\sqrt{d-1}$, $\rho=\sigma^{-2}<1$, and substitute $x=2\sigma\cos\theta$, $0\le\theta\le\pi$. Expanding the denominator gives
\[
d^2-4\sigma^2\cos^2\theta
=\sigma^4|1-\rho e^{2i\theta}|^2,
\]
because $d=\sigma^2+1$. Including the Jacobian transforms the density into
\[
\frac{2d}{\pi(d-1)}
\frac{\sin^2\theta}{|1-\rho e^{2i\theta}|^2}\dd\theta.
\]

Define $U_j$ by $U_j(\cos\theta)=\sin((j+1)\theta)/\sin\theta$ and $U_{-1}=0$. The sine addition formula gives $U_{j+1}(z)=2zU_j(z)-U_{j-1}(z)$, so these are polynomials. Checking $k=1,2$ and then using this recurrence proves
\[
q_k(2\sigma\cos\theta)=\sigma^k(U_k-\rho U_{k-2})(\cos\theta),\quad k\ge1.
\]
Consequently
\[
\sin\theta\,q_k(2\sigma\cos\theta)
=\sigma^k\operatorname{Im}\left(e^{i(k-1)\theta}(e^{2i\theta}-\rho)\right).
\]
The quotient needed for integration has a uniformly convergent geometric expansion:
\[
\frac{e^{2i\theta}-\rho}{|1-\rho e^{2i\theta}|^2}
=\frac1{e^{-2i\theta}-\rho}
=\sum_{j\ge0}\rho^je^{2i(j+1)\theta}.
\]
Termwise integration is therefore valid, and for $k\ge1$ it gives
\[
\int q_k\dd\nu_d=
\frac{2d\sigma^k}{\pi(d-1)}
\sum_{j\ge0}\rho^j\int_0^\pi\sin\theta\sin((k+1+2j)\theta)\dd\theta=0.
\]
Each sine frequency in the second factor is at least $2$, orthogonal to $\sin\theta$ on $[0,\pi]$.

For the mass, multiply the two geometric series to obtain
\[
\frac1{|1-\rho e^{2i\theta}|^2}
=\frac1{1-\rho^2}\sum_{j\in\Z}\rho^{|j|}e^{2ij\theta}.
\]
Uniform convergence again allows integration. Since $\sin^2\theta=(1-\cos2\theta)/2$, only frequencies $0,2,-2$ contribute, giving integral $\pi/(2(1+\rho))$. Multiplying by $2d/(\pi(d-1))$ yields $1$, because $1+\rho=d/(d-1)$. The density is non-negative with mass one, and its polynomial integrals agree with $L$, as required.
\end{proof}

This is the \emph{Kesten--McKay law}. Its density is positive in the interior of $[-R,R]$, where $R=2\sqrt{d-1}$. Thus its even moments have root growth $a_{2r}^{1/(2r)}\to R$: the upper bound is $R$, and any interval near $R$ with positive mass gives the matching lower limit. Exercise~\ref{ex:v8c07c} writes out that argument. The return probabilities $p_{2r}=a_{2r}/d^{2r}$ consequently have root limit $R/d$.

This return-growth description is one meaning of the tree's spectral radius. Chapter~\ref{ch:operators} develops the operator language behind the name. For the line $T_2$, the return count is $\binom{2r}{r}$, whose $2r$th root tends to $2$, giving the same formula at degree two. For $T_{2k}$, the walk also describes the free group on $k$ generators; Volume~12, Chapter~8 gives its generating-function proof. Rescaling $\nu_d$ by $\sqrt{d-1}$ gives the semicircle limit in Exercise~\ref{ex:kmsemicircle}, connected with Volume~12, Chapter~9.

\begin{theorem}[McKay, 1981]\label{thm:mckay}
Fix $d\ge3$. Suppose $G_N$ is a sequence of $d$-regular graphs on $N\to\infty$ vertices such that, for every fixed radius $r$, the fraction of vertices whose radius-$r$ ball contains a cycle tends to zero. Then the empirical spectral measures
$N^{-1}\sum_i\delta_{\lambda_i(A(G_N))}$ converge weakly to $\nu_d$.
\end{theorem}
\begin{proof}
Fix a non-negative integer $k$ and choose $r\ge k$. At a vertex with a cycle-free radius-$r$ ball, that ball agrees with the tree up to depth $r$: every vertex below the boundary has all $d$ neighbours inside the ball. A closed walk of length $k$ stays at distance at most $k/2$ from its start. Its count there is therefore $a_k$.
At any other vertex the count is between zero and $d^k$; also $a_k\le d^k$. Taking the trace gives the bound
\[
\left|\frac1N\sum_i\lambda_i^k-a_k\right|
\le2d^k\frac{\#\{x:B(x,r)\text{ contains a cycle}\}}N\longrightarrow0.
\]
All eigenvalues lie in $[-d,d]$ by the row-sum bound, and $\nu_d$ is supported there too. Polynomial integrals converge by the moment calculation. Uniform polynomial approximation on this interval, as proved by Weierstrass, then gives convergence of integrals of every continuous function there. This is weak convergence, since all measures share this compact support.
\end{proof}

For uniformly random regular graphs, fixed-length cycle counts remain bounded in probability; more precisely, each length-$\ell$ count converges to a Poisson law of mean $(d-1)^\ell/(2\ell)$, a result of Bollob\'as. A fixed-radius cyclic ball lies within bounded distance of a cycle of bounded length. Thus the fraction of such balls tends to zero in probability, giving the Kesten--McKay limit for random graphs. A histogram illustrates the limit but does not prove it.

\begin{corollary}\label{cor:serre}
Under Theorem~\ref{thm:mckay}'s hypothesis, for each $\varepsilon>0$ there is $c(\varepsilon,d)>0$ such that, for all sufficiently large $N$, at least that fraction of the eigenvalues is at least $2\sqrt{d-1}-\varepsilon$.
\end{corollary}
\begin{proof}
Let $R=2\sqrt{d-1}$. Choose a continuous function $g$ taking values in $[0,1]$, zero on $(-\infty,R-\varepsilon]$ and one on $[R-\varepsilon/2,\infty)$. The positive density near $R$ gives $a=\int g\dd\nu_d>0$. Weak convergence makes the empirical integral at least $a/2$ eventually. Since $g$ vanishes below the threshold and never exceeds one, that integral is bounded by the desired fraction. Take $c=a/2$.
\end{proof}

The top edge of the tree's spectrum attracts a positive fraction of the eigenvalues under this local-tree hypothesis. Serre proved a broader result for all fixed-degree regular graphs; the present corollary proves the version needed here. A Ramanujan graph allows no nontrivial eigenvalue beyond either tree edge, a stronger requirement than convergence of the empirical measure.

\begin{figure}[ht]
\centering
\fig{ch09-kestenmckay}
\caption{Left: histogram of the $3000$ eigenvalues of a random $4$-regular graph and the Kesten--McKay density $\nu_4$ (rust), supported on $[-2\sqrt3,2\sqrt3]$; the trivial eigenvalue $4$ is the isolated dot. Right: the
poles $u$ of the Ihara zeta function of a random $3$-regular graph on $300$ vertices (Section~\ref{sec:ihara}). The trivial poles $u=1$ and $u=\frac12$ (coral) lie on the real axis, and so do the poles $\pm1$ of the factor $(1-u^2)^{m-n}$ (not drawn); all the others lie on the circle
$|u|=\frac1{\sqrt2}$: this graph is Ramanujan, its zeta function satisfies the Riemann hypothesis.}
\label{fig:kestenmckay}
\end{figure}

\section{Ihara's zeta function}\label{sec:ihara}

Ihara introduced his zeta function in 1966 while studying groups over $p$-adic fields. It counts closed geodesics. On a graph, a geodesic is represented by a walk without immediate reversals, including at the join between its last and first edges. Serre, Sunada, Hashimoto and Bass developed the graph formulation and its determinant identities. We can prove the identities using finite matrices.

For this section let $G$ be a finite connected simple graph, with $n$ vertices and $m$ edges. Each edge supplies two orientations. Write $o(e)$ and $t(e)$ for the origin and terminus of an oriented edge, and $\bar e$ for its reverse. Define the $2m\times2m$ non-backtracking matrix by
\[
B(e,f)=\begin{cases}1,&t(e)=o(f),\ f\ne\bar e,\\0,&\text{otherwise}.\end{cases}
\]
A cycle of length $k$ is a sequence $(e_1,\ldots,e_k)$ with $B(e_i,e_{i+1})=1$, using cyclic indices. These are rooted oriented cycles, counted by $N_k=\operatorname{tr}B^k$. Identify sequences differing by a cyclic rotation, but keep the two directions distinct. A class is \emph{primitive} if it is not a repeated traversal of a shorter cycle. Let $\pi(k)$ count primitive classes of length $k$. The Euler product is
\[
Z_G(u)=\prod_{[P]}(1-u^{\ell(P)})^{-1}
=\prod_{k\ge1}(1-u^k)^{-\pi(k)},
\]
with one factor for each primitive class.

\begin{proposition}\label{prop:iharadet}
In a neighbourhood of $u=0$,
\[
\log Z_G(u)=\sum_{k\ge1}\frac{N_k}{k}u^k,\qquad
Z_G(u)=\det(I-uB)^{-1}.
\]
In particular its continuation is the reciprocal of a polynomial.
\end{proposition}
\begin{proof}
Each rooted cycle has a least period $\ell$ dividing its length $k$, found from the subgroup of rotations fixing its edge sequence. Its first $\ell$ edges form a primitive cycle. A primitive class of length $\ell$ gives exactly $\ell$ distinct rooted cycles whenever $\ell\mid k$, not $k$ distinct ones. Hence
\begin{equation}\label{eq:nkpi}
N_k=\sum_{\ell\mid k}\ell\pi(\ell).
\end{equation}
Let $b$ be the maximum row sum of $B$, taking $b=0$ if there are no oriented edges. We have $N_k\le2m b^k$. For $|u|<1/\max(1,b)$, absolute convergence justifies expanding and regrouping the logarithms of the product:
\[
\log Z_G(u)=\sum_{\ell\ge1}\pi(\ell)\sum_{j\ge1}\frac{u^{j\ell}}j
=\sum_{k\ge1}\frac{u^k}{k}\sum_{\ell\mid k}\ell\pi(\ell).
\]
Triangularise $B$ over $\C$, using the argument preceding Theorem~\ref{linear:thm:spectral}, and call its eigenvalues $\eta_i$. Triangular matrices give $\operatorname{tr}B^k=\sum_i\eta_i^k$ and $\det(I-uB)=\prod_i(1-u\eta_i)$. Expanding each logarithm near zero yields
$\sum_k\operatorname{tr}(B^k)u^k/k=-\log\det(I-uB)$.
Exponentiation proves the identity.
\end{proof}

\begin{theorem}[Ihara, 1966; Bass, 1992]\label{thm:bass}
For a $(q+1)$-regular graph,
\[
\det(I-uB)=(1-u^2)^{m-n}\det(I-uA+qu^2I).
\]
If $m<n$, interpret the right side as a rational expression with its removable factors cancelled; the resulting identity holds at every $u$.
\end{theorem}
\begin{proof}
Let $S,T$ be the $n\times2m$ origin and terminus incidence matrices, and let $R$ reverse oriented edges. Checking their entries gives
\[
ST^{\mathsf T}=A,\qquad TT^{\mathsf T}=(q+1)I,
\qquad SR=T,\qquad B=T^{\mathsf T}S-R.
\]
Put $M=I+uR$. Each pair of orientations contributes a block $\begin{pmatrix}1&u\\u&1\end{pmatrix}$, so $\det M=(1-u^2)^m$ and $M^{-1}=(I-uR)/(1-u^2)$ for $u^2\ne1$.
We also need $\det(I-XY)=\det(I-YX)$ for rectangular compatible $X,Y$. Eliminate either diagonal block in $\begin{pmatrix}I&X\\Y&I\end{pmatrix}$; the two resulting determinants prove this identity.
Taking $W=SM^{-1}T^{\mathsf T}$, factor and apply it to get
\[
\begin{aligned}
\det(I-uB)&=\det(M-uT^{\mathsf T}S)\\
&=(1-u^2)^m\det(I_n-uW).
\end{aligned}
\]
The incidence identities give $W=(A-u(q+1)I)/(1-u^2)$. Multiplying the last determinant by $(1-u^2)^n$ therefore yields $\det(I-uA+qu^2I)$. This proves the formula away from $u=\pm1$. If $m\ge n$, both sides are polynomials, so the identity holds everywhere. If $m<n$, the equality away from those two points shows that all apparent singularities of the rational expression are removable, with its extension equal to the polynomial on the left.
\end{proof}

For degree $q+1\ge2$, we have $m\ge n$. Each adjacency eigenvalue $\lambda_j$ contributes a factor
$1-\lambda_ju+qu^2=(1-\mu_ju)(1-\mu'_ju)$,
where $\mu_j,\mu'_j$ solve $z^2-\lambda_jz+q=0$. Comparing characteristic factors gives the eigenvalues of $B$: these pairs and $m-n$ additional copies of each of $1,-1$. Thus
\begin{equation}\label{eq:nkspec}
N_k=\sum_{j=1}^n(\mu_j^k+\mu_j'^k)+(m-n)(1+(-1)^k).
\end{equation}
The degree restriction here prevents a negative number of additional eigenvalues; the determinant identity itself also covered the lower-degree cases.

\section{The Riemann hypothesis for graphs}

Take a connected $(q+1)$-regular graph with $q\ge2$. It is \emph{Ramanujan} when every adjacency eigenvalue apart from $q+1$ and, in the bipartite case, $-(q+1)$ lies in $[-2\sqrt q,2\sqrt q]$. The trivial positive eigenvalue gives factor
$1-(q+1)u+qu^2=(1-u)(1-qu)$,
so its poles occur at $u=1,1/q$. Under $u=q^{-s}$ these correspond to $s=0,1$, with repetitions in the imaginary direction. The circle $|u|=q^{-1/2}$ corresponds to the vertical line $\operatorname{Re}s=1/2$.

\portrait{riemann}{Bernhard Riemann (1826--1866). Photograph, 1863. Public domain.}

\begin{theorem}[Sunada, 1986]\label{thm:graphRH}
A connected $(q+1)$-regular graph with $q\ge2$ is Ramanujan if and only if every pole of $Z_G(u)$ in the annulus $1/q<|u|<1$ lies on $|u|=q^{-1/2}$. Equivalently, after putting $u=q^{-s}$, all poles in $0<\operatorname{Re}s<1$ lie on $\operatorname{Re}s=1/2$.
\end{theorem}
\begin{proof}
Bass's formula makes $Z_G$ the reciprocal of a polynomial with factors $1-u^2$ and $1-\lambda u+qu^2$. There are no numerator zeros to cancel their roots. The trivial eigenvalues $\lambda=\pm(q+1)$ give roots of modulus $1$ or $1/q$, on the annulus boundaries.
For $|\lambda|\le2\sqrt q$, the quadratic roots are conjugate, or a repeated real root, with product $1/q$. Both have modulus $1/\sqrt q$.
If $2\sqrt q<|\lambda|<q+1$, replace $u$ by $-u$ when $\lambda<0$. The polynomial becomes $qu^2-|\lambda|u+1$. It has two distinct positive roots. It is positive at $1/q$ and at $1$, with respective values $(q+1-|\lambda|)/q$ and $q+1-|\lambda|$. Its vertex lies strictly between these points and has negative value. Both roots therefore lie inside $(1/q,1)$, and their product $1/q$ puts one on either side of $1/\sqrt q$. They produce poles off the critical circle. These cases exhaust all eigenvalues by the row-sum bound, proving both directions.
\end{proof}

For the ordinary Riemann zeta function, a square-root-scale error in prime counting is tied to its Riemann hypothesis (Volume~6, Theorem~7.10). The graph analogue follows from the finite spectral formula and the divisor identity for primitive cycles.

\begin{theorem}[prime number theorem for graphs]\label{thm:graphpnt}
Let $G$ be a connected nonbipartite Ramanujan $(q+1)$-regular graph on $n$ vertices, with $q\ge2$. For every integer $k\ge2$,
\[
\left|\pi(k)-\frac{q^k}{k}\right|\le\frac{6nq^{k/2}}{k}.
\]
\end{theorem}
\begin{proof}
In~\eqref{eq:nkspec}, the top adjacency eigenvalue contributes $q^k+1$. All other pairs have $|\mu_j|=|\mu'_j|=\sqrt q$, since there is no eigenvalue $-(q+1)$. Also $m-n=n(q-1)/2$. For $k\ge2$ these observations give
\[
|N_k-q^k|\le1+2(n-1)q^{k/2}+n(q-1)\le3nq^{k/2}.
\]
Equation~\eqref{eq:nkpi} says that $k\pi(k)$ is $N_k$ minus the contributions of proper divisors. That subtraction is non-negative and at most $\sum_{\ell\le k/2}N_\ell$. A rooted non-backtracking cycle of length $\ell$ chooses its first oriented edge in $2m$ ways and its remaining steps in at most $q^{\ell-1}$ ways. Therefore
\[
\sum_{\ell\le k/2}N_\ell
\le\frac{2m}{q-1}q^{k/2}
=\frac{n(q+1)}{q-1}q^{k/2}\le3nq^{k/2}.
\]
Combining the two errors and dividing by $k$ proves the bound.
\end{proof}

\begin{beauty}
A primitive cycle supplies an Euler-product factor, and a matrix eigenvalue supplies a quadratic determinant factor. Bass's identity makes them descriptions of the same function. Its critical circle is then a quadratic-root calculation, while its primitive-cycle count follows by separating the largest spectral term from the rest. The analogy with primes has an exact formula behind it.
\end{beauty}

\section{Ramanujan graphs}

Alon--Boppana forbids a fixed improvement below $2\sqrt q$ along a growing family. Ramanujan graphs reach that bound for every nontrivial eigenvalue. The results below describe ways to obtain them. Their construction proofs use additional mathematics and are not premises for any proof in this chapter.

A Cayley graph turns a generating set into edges: each group element is joined to its products with the generators. This connects a finite network to the algebra of the group. Cayley's name belongs here as well as to matrix identities.

\portrait{cayley}{Arthur Cayley (1821--1895). Photograph by Herbert Rose Barraud, before 1883. Public domain, via Wikimedia Commons.}

\begin{remark}[explicit constructions]
Fix a prime $p\equiv1\pmod4$, and choose a distinct sufficiently large prime $\ell\equiv1\pmod4$. The LPS construction begins with the $p+1$ integer solutions of
$a_0^2+a_1^2+a_2^2+a_3^2=p$
for which $a_0>0$ is odd and $a_1,a_2,a_3$ are even. Jacobi's four-square formula gives this count. A quaternion representation sends these solutions to projective matrices over $\mathbf F_\ell$, closed under inverses. Joining group elements by multiplication with these generators defines the associated Cayley graph, of degree $p+1$. Choosing $\ell$ sufficiently large avoids coincidences of generators or loops. The relevant group is $\mathrm{PSL}_2$ or $\mathrm{PGL}_2$, according to the quadratic-residue case.

The nontrivial graph eigenvalues are related to Hecke eigenvalues of weight-two modular forms, giving the bound $2\sqrt p$. Eichler's theorem, using Weil's theorem for curves, supplies the case needed here; Deligne proved the broader Ramanujan conjecture. The bipartite LPS families have girth at least $(4/3)\log_p N$ up to an additive constant. Margulis obtained such families independently, and Morgenstern extended the construction to prime-power $q$ and degree $q+1$. Davidoff, Sarnak and Valette's book in the sources develops the LPS construction with its hypotheses in full.
\end{remark}

\begin{remark}[existence for all degrees, and random graphs]
Marcus, Spielman and Srivastava's 2013 work gives bipartite Ramanujan families in every degree $d\ge3$, through signings and interlacing families of polynomials. Their method also led to the solution of Kadison--Singer. Exercise~\ref{ex:v8c07d} explains the elementary spectral relation between a signing and its double cover; finding an appropriate signing is the deeper step.
Friedman's theorem gives near-Ramanujan bounds with probability tending to one for fixed degree. Huang, McKenzie and Yau's 2024 preprint identifies the edge fluctuations at scale $N^{-2/3}$, with degree-dependent normalisation, and proves their limiting Tracy--Widom laws and asymptotic independence. These yield the approximately $69\%$ limiting Ramanujan probability quoted at the start, for uniformly random simple graphs and each fixed $d\ge3$.
\end{remark}

\begin{remark}[Cayley graphs of $\mathrm{SL}_2(\mathbf F_p)$]
For odd primes $p$, use the four generators
\[
S=\left\{\begin{pmatrix}1&\pm1\\0&1\end{pmatrix},
\begin{pmatrix}1&0\\\pm1&1\end{pmatrix}\right\}.
\]
They give connected $4$-regular graphs on $p(p^2-1)$ vertices, as Exercise~\ref{ex:sl2gen} proves. These graphs have a spectral gap uniform in $p$, obtained from Selberg's $3/16$ theorem for the modular surfaces by a comparison with the finite graphs. Lubotzky's book explains the transfer. The generators merge in pairs at $p=2$, so the simple graph then has degree $2$, not $4$.
Bourgain and Gamburd proved uniform expansion for reductions of a fixed finite symmetric integer generating set whose subgroup is Zariski dense in $\mathrm{SL}_2$, for sufficiently large primes. Their argument uses additive combinatorics. Uniform expansion does not assert the sharp Ramanujan bound; when using a one-sided gap for walks, laziness supplies convergence as in Exercise~\ref{ex:v8c06d}.
\end{remark}

\section{Beyond finite spectra}

Sparse graphs have supplied several different uses for a gap: mixing a walk, avoiding a set of bad seeds, and concentrating sample averages. Return counts on the tree then fixed the scale of a possible gap. Ihara's product and Bass's determinant formula connected that scale to poles and primitive cycles. The construction results go further than these finite-matrix proofs, and have been stated with their sources rather than used without proof.

The next chapter leaves finite matrices for bounded operators. In infinite dimension the spectrum can contain values that are not eigenvalues at all. We will need to distinguish failure of invertibility from the existence of an eigenvector before carrying the spectral story onward.

\begin{figure}[ht]
\centering\fig{ch07-rw-returns}
\caption{A length-four return walk on a tree has distance pattern $0,1,0,1,0$ or $0,1,2,1,0$. The patterns contribute $d^2$ and $d(d-1)$ walks, respectively, giving $a_4=d(2d-1)$.}
\label{fig:v8c07returns}
\end{figure}
\begin{figure}[ht]
\centering\fig{ch07-rw-cover}
\caption{A negative sign on one triangle edge makes that edge cross between the two sheets of a double cover. The resulting graph is a six-cycle. Sheet-symmetric vectors retain the original spectrum; sheet-antisymmetric vectors see the signed matrix.}
\label{fig:v8c07cover}
\end{figure}

\section{Exercises}

\begin{exercise}[$\star$]\label{ex:bipartiteeig}
For a connected $d$-regular graph with $d>0$, prove that $-d$ occurs in the adjacency spectrum exactly when the graph is bipartite. Show that a usual walk started at a vertex then fails to converge to uniform. For the lazy transition $(I+P)/2$, locate every nonconstant eigenvalue in $[0,(1+\lambda_2/d)/2]$.
\end{exercise}
\begin{exercise}[$\star\star$]\label{ex:amplify}
A $k$-bit randomised algorithm always answers ``no'' on a no instance, and answers ``yes'' on at least half its seeds on a yes instance. Put the $2^k$ seeds on a $d$-regular graph with $\alpha\le1/4$. When $d$ is a power of two, design an amplification using $k+t\log_2d$ random bits and error at most $(3/4)^t$. For other degrees, explain the obstruction to a fixed exact bit budget for uniform neighbour choices, and give a uniform rejection sampler with expected fewer than $2\lceil\log_2d\rceil$ bits per step.
\end{exercise}
\begin{exercise}[$\star$]\label{ex:treemoments}
Count the tree's length-four return walks to obtain $d(2d-1)$. Recover the count from $\int q_4\dd\nu_d=0$, using the polynomial recurrence. Continue to length six and prove $a_6=d(5d^2-6d+2)$.
\end{exercise}
\begin{exercise}[$\star\star$]\label{ex:kmsemicircle}
Push $\nu_d$ forward by $x\mapsto x/\sqrt{d-1}$. Prove convergence to the probability density $(2\pi)^{-1}\sqrt{4-y^2}\mathbf1_{|y|\le2}$ as $d\to\infty$. Show that the normalised even moments $a_{2k}/(d-1)^k$ tend to $C_k=(k+1)^{-1}\binom{2k}{k}$.
\end{exercise}
\begin{exercise}[$\star$]\label{ex:cycleihara}
Find $Z_{C_n}(u)$ both from primitive cycles and from Bass's formula at $q=1$. Both methods should give $(1-u^n)^{-2}$.
\end{exercise}
\begin{exercise}[$\star\star$]\label{ex:k4}
Verify that $K_4$ is Ramanujan. Compute its non-backtracking spectrum and the counts $N_3=N_4=24$. Draw the primitive cycles to verify $\pi(3)=8$ and $\pi(4)=6$.
\end{exercise}
\begin{exercise}[$\star\star$]\label{ex:irregularbass}
For any finite connected simple graph, put $Q=\operatorname{diag}(\deg x-1)$. Prove
$\det(I-uB)=(1-u^2)^{m-n}\det(I-uA+u^2Q)$,
interpreting removable factors as in Theorem~\ref{thm:bass} when $m<n$.
\end{exercise}
\begin{exercise}[$\star\star$]\label{ex:sl2gen}
For every prime $p$, count $\mathrm{SL}_2(\mathbf F_p)$ and show that $U=\begin{pmatrix}1&1\\0&1\end{pmatrix}$ and $L=\begin{pmatrix}1&0\\1&1\end{pmatrix}$ generate it. Deduce connectedness of the Cayley graph with $U^{\pm1},L^{\pm1}$, and degree four for $p\ge3$.
\end{exercise}

\begin{exercise}[$\star\star$]\label{ex:v8c07a}
For a usual walk on $K_4$ started uniformly, let $B$ contain two fixed vertices. Compute the exact probability of staying in $B$ from time $0$ to time $t$. Compare it with Proposition~\ref{prop:hitting}'s bound using $\beta=1/2$ and $\alpha=1/3$.
\end{exercise}
\begin{exercise}[$\star\star$]\label{ex:v8c07b}
Show that $K_{3,3}$ is Ramanujan. Determine its non-backtracking eigenvalues and compute $N_4$ and $\pi(4)$. Explain why all its odd cycle counts vanish.
\end{exercise}
\begin{exercise}[$\star\star$]\label{ex:v8c07c}
Let a probability measure be supported on $[-R,R]$, with $R>0$, and give positive mass to $(R-\delta,R]$ for every $0<\delta<R$. Prove that its even moments $m_{2r}$ satisfy $m_{2r}^{1/(2r)}\to R$. Apply this to the tree return probabilities, and give the resulting limit for degree $3$.
\end{exercise}
\begin{exercise}[$\star\star$]\label{ex:v8c07d}
Sign each edge of a graph by $+1$ or $-1$, and write $A_s$ for its signed adjacency matrix. Form a double cover on $(v,0),(v,1)$: a positive edge joins equal sheet indices, and a negative edge joins opposite indices. Prove that the cover's spectrum, with multiplicities, is the union of the spectra of $A$ and $A_s$. Work out a triangle with one negative edge and identify its cover.
\end{exercise}

\section{Hints}

\begin{hint}{ex:bipartiteeig}
Use the connected-graph equivalence already proved in Exercise~\ref{linear:ex:bipartite}. Apply the affine map $x\mapsto(1+x/d)/2$ to each eigenvalue for the lazy walk.
\end{hint}
\begin{hint}{ex:amplify}
Answer yes if any visited seed does. For a general degree, sample a binary integer in $[0,2^b)$ with $b=\lceil\log_2d\rceil$ and retry if it is at least $d$.
\end{hint}
\begin{hint}{ex:treemoments}
There are two length-four distance patterns, shown in Figure~\ref{fig:v8c07returns}. Generate $q_4$ and $q_6$ recursively and use the vanishing integrals.
\end{hint}
\begin{hint}{ex:kmsemicircle}
After substituting $x=\sqrt{d-1}\,y$, separate the semicircle density from its rational multiplier. For its moments, use $y=2\cos\theta$ and the Wallis recurrence.
\end{hint}
\begin{hint}{ex:cycleihara}
An oriented non-backtracking walk on a cycle cannot change direction. In the determinant calculation factor each quadratic using the roots of unity, without complex conjugating $u$.
\end{hint}
\begin{hint}{ex:k4}
The three nontrivial adjacency eigenvalues are $-1$. The associated non-backtracking roots satisfy $z^2+z+2=0$.
\end{hint}
\begin{hint}{ex:irregularbass}
In the incidence calculation replace $(q+1)I$ by the degree matrix $Q+I$ and keep the other identities unchanged.
\end{hint}
\begin{hint}{ex:sl2gen}
Choose a non-zero first column, then solve one nonzero linear equation for the second. Powers of $U,L$ implement elementary row additions.
\end{hint}
\begin{hint}{ex:v8c07a}
From a vertex in $B$, exactly one of its three neighbours is still in $B$.
\end{hint}
\begin{hint}{ex:v8c07b}
The adjacency spectrum is $3,-3,0,0,0,0$. Use the quadratic pairs at $q=2$ and remember the extra $\pm1$ factors.
\end{hint}
\begin{hint}{ex:v8c07c}
Bound the moment above by $R^{2r}$ and below by $(R-\delta)^{2r}$ times a fixed positive mass.
\end{hint}
\begin{hint}{ex:v8c07d}
Write the cover matrix using the positive-edge adjacency $A_+$ and the negative-edge adjacency $A_-$. Test vectors of the forms $(x,x)$ and $(x,-x)$.
\end{hint}

\section{Solutions}

\begin{solution}{ex:bipartiteeig}
\idea The negative extreme eigenvector identifies the two sides; laziness moves that eigenvalue to zero.
\writeup Exercise~\ref{linear:ex:bipartite} proves that a connected graph has $-\lambda_1$ in its spectrum exactly when it is bipartite. Here $\lambda_1=d$. For a point start in a bipartite graph, all mass lies on one side at each time, alternating sides. Regularity makes the two sides equally large, because each has $|E|/d$ vertices. Thus the total variation distance from uniform is at least $1/2$ at every step. A uniform initial law is already stationary; the failure assertion concerns the specified point start.
The lazy eigenvalues are $(1+\lambda_i/d)/2$. Since $-d\le\lambda_i\le d$ and the constant eigenvalue $d$ is simple, every nonconstant one lies between zero and $(1+\lambda_2/d)/2<1$. Hence the lazy walk converges from every starting law.
\end{solution}

\begin{solution}{ex:amplify}
\idea A one-sided procedure fails only if the walk never visits a successful seed.
\writeup Choose $X_0$ uniformly, take $t$ uniform-neighbour steps, and run the algorithm at each visited seed. Answer yes if any run does, otherwise no. On a no instance there is no erroneous yes. On a yes instance the bad set has $\beta\le1/2$, so the hitting estimate gives error at most $\beta(\beta+\alpha)^t\le\tfrac12(3/4)^t$.
If $d=2^b$, the initial seed costs $k$ bits and each step costs exactly $b$, giving the requested budget. If $d$ has an odd factor, a uniform choice among $d$ outcomes cannot be implemented from any fixed finite number $b$ of fair bits: each outcome probability is an integer multiple of $2^{-b}$, whereas $1/d$ is not. In particular, rounding $\log_2d$ upward does not solve this exact sampling problem.
For general non-power-of-two $d$, draw a uniform $b$-bit integer with $b=\lceil\log_2d\rceil$; accept values $0,\ldots,d-1$ and otherwise retry independently. Conditional on acceptance all $d$ outcomes are equally likely. The acceptance probability is $d/2^b>1/2$, so the expected number of trials is $2^b/d<2$. The expected bits per step are therefore $b2^b/d<2b$. Sampling terminates almost surely and preserves the same error estimate, with this expected rather than fixed bit budget.
\end{solution}

\begin{solution}{ex:treemoments}
\idea Count distance patterns, then extract moments from monic polynomial identities.
\writeup A four-step return either returns after two steps and leaves again, in $d^2$ ways, or travels to distance two and returns along the same path, in $d(d-1)$ ways. Their sum is $2d^2-d$.
The recurrence gives
\[
q_3=x^3-(2d-1)x,\qquad
q_4=x^4-(3d-2)x^2+d(d-1).
\]
With $a_0=1$, $a_2=d$ and $L(q_4)=0$, we recover $a_4=(3d-2)d-d(d-1)=2d^2-d$. Continuing the recurrence gives
\[
\begin{aligned}
q_5&=x^5-(4d-3)x^3+(3d^2-4d+1)x,\\
q_6&=x^6-(5d-4)x^4+(6d^2-9d+3)x^2-d(d-1)^2.
\end{aligned}
\]
Applying $L(q_6)=0$ gives
\[
a_6=(5d-4)(2d^2-d)-(6d^2-9d+3)d+d(d-1)^2
=d(5d^2-6d+2).
\]
At degree three the counts are $a_4=15$ and $a_6=87$, obtained here by exact identities.
\end{solution}

\begin{solution}{ex:kmsemicircle}
\idea The rescaled density differs from a semicircle by a multiplier tending uniformly to one.
\writeup The substitution gives density
\[
\frac{d(d-1)}{d^2-(d-1)y^2}\,
\frac{\sqrt{4-y^2}}{2\pi}\mathbf1_{|y|\le2}.
\]
On $[-2,2]$ the denominator is at least $(d-2)^2$, and the rational multiplier tends uniformly to $1$. Thus densities, integrals of bounded continuous functions, and every fixed moment converge on this compact interval.
For the limiting even moment, write $J_k=\int_0^\pi\cos^{2k}\theta\dd\theta$. Integration by parts gives $J_{k+1}=(2k+1)J_k/(2k+2)$ and $J_0=\pi$, so $J_k/\pi=4^{-k}\binom{2k}{k}$. The substitution $y=2\cos\theta$ gives
\[
\frac1{2\pi}\int_{-2}^2y^{2k}\sqrt{4-y^2}\dd y
=\frac{2\cdot4^k}{\pi}(J_k-J_{k+1})
=\frac1{k+1}\binom{2k}{k}.
\]
The rescaled moment is $a_{2k}/(d-1)^k$, proving the desired Catalan limit. The counting counterpart is that each distance pattern of a returning tree walk is a Dyck path; its $k$ outward choices have leading count $d^k$ as $d$ grows, with $k$ fixed.
\end{solution}

\begin{solution}{ex:cycleihara}
\idea Each direction around the cycle supplies one primitive class.
\writeup A non-backtracking walk on $C_n$ must continue in its chosen direction. A closed one traverses the full cycle some positive number of times. There are therefore exactly two primitive classes, each of length $n$, giving $Z=(1-u^n)^{-2}$.
For the determinant proof, $q=1$ and $m=n$. Put $\omega=e^{2\pi i/n}$. The adjacency eigenvalues give
\[
\begin{aligned}
\det(I-uB)&=\prod_{j=0}^{n-1}(1-u\omega^j)(1-u\omega^{-j})\\
&=(1-u^n)^2.
\end{aligned}
\]
Each of the two products is $1-u^n$, by factorisation over the $n$th roots of unity. This polynomial calculation is valid for complex $u$ as well as real $u$.
\end{solution}

\begin{solution}{ex:k4}
\idea Each adjacency eigenvalue supplies two non-backtracking eigenvalues through Bass's quadratic.
\writeup The adjacency spectrum is $3,-1,-1,-1$. Since $1\le2\sqrt2$, the graph is Ramanujan. Here $q=2$, $n=4$, $m=6$. The eigenvalue $3$ supplies roots $2,1$, and each $-1$ supplies roots $z,\bar z$ of $z^2+z+2=0$. Add two further copies of each of $1,-1$ because $m-n=2$. This lists all $12$ eigenvalues with multiplicity.
The quadratic gives $z^3=2-z$ and $z^4=3z+2$. Since $z+\bar z=-1$, we get $z^3+\bar z^3=5$ and $z^4+\bar z^4=1$. Taking traces gives $N_3=8+1+3\cdot5=24$ and $N_4=16+1+3\cdot1+4=24$. A simple graph has no non-backtracking cycles of length one or two, so the divisor identity gives $\pi(3)=24/3=8$ and $\pi(4)=24/4=6$. Directly, the four triangles have two orientations each, and the three undirected Hamiltonian four-cycles have two orientations each.
\end{solution}

\begin{solution}{ex:irregularbass}
\idea Keep the degree matrix in the incidence proof instead of replacing it by a scalar.
\writeup With the same $S,T,R,M$, we now have $TT^{\mathsf T}=Q+I$. All other incidence identities remain valid, giving
\[
SM^{-1}T^{\mathsf T}=\frac{A-u(Q+I)}{1-u^2}.
\]
Consequently the numerator matrix for $I-uSM^{-1}T^{\mathsf T}$ is
$(1-u^2)I-uA+u^2(Q+I)=I-uA+u^2Q$.
Multiplying its determinant by $\det M=(1-u^2)^m$ supplies the claimed factor $(1-u^2)^{m-n}$. The equality first holds for $u^2\ne1$; the left polynomial identifies the removable extension at the remaining points, including trees where $m-n=-1$.
\end{solution}

\begin{solution}{ex:sl2gen}
\idea Every determinant-one matrix can be reduced to the identity by the two kinds of row addition.
\writeup The first column can be any non-zero vector, giving $p^2-1$ choices. For a fixed first column, the equation making the determinant $1$ is a nonzero linear equation in the second column's two coordinates, with $p$ solutions. Thus the group order is $p(p^2-1)$.
Let $g=\begin{pmatrix}a&b\\c&d\end{pmatrix}$ have determinant one. Left multiplication by $U^x$ adds $x$ times row two to row one; multiplication by $L^y$ adds $y$ times row one to row two. Powers realise every field coefficient because the field is $\mathbf F_p$. If $c\ne0$, choose $x$ so that $a+xc=1$. Next use a lower row addition to make the lower-left entry zero. Determinant one then forces the lower-right entry to be one. A final upper row addition removes the upper-right entry. If initially $c=0$, then $a\ne0$ and multiplication by $L$ first makes the lower-left entry non-zero. We have reduced $g$ to $I$, proving generation for every prime. The undirected Cayley graph is therefore connected. For $p\ge3$ the four elements $U,U^{-1},L,L^{-1}$ are distinct and nonidentity, giving degree four.
\end{solution}

\begin{solution}{ex:v8c07a}
\idea Staying inside the chosen pair has the same conditional probability at every step.
\writeup The uniform start lands in $B$ with probability $1/2$. From either vertex of $B$, just one of the three neighbours belongs to $B$, so the conditional probability of staying for the next step is $1/3$. The exact path-sum probability is $(1/2)(1/3)^t$. The proposition gives $(1/2)(1/2+1/3)^t=(1/2)(5/6)^t$. It is a valid but looser upper bound; the compression in this example has norm $1/3$, whereas the triangle estimate used in the proposition only bounds it by $5/6$.
\end{solution}

\begin{solution}{ex:v8c07b}
\idea Keep both trivial adjacency eigenvalues; bipartiteness contributes an additional extreme pair.
\writeup Exercise~\ref{ex:v8c05a} gives adjacency spectrum $3,-3,0,0,0,0$. The nontrivial zeros lie within the Ramanujan interval for $q=2$. Bass's pairs are $2,1$ from $3$, $-2,-1$ from $-3$, and $i\sqrt2,-i\sqrt2$ each four times from the zeros. Since $m=9,n=6$, add three copies each of $1,-1$. This totals $18$ eigenvalues.
Their fourth powers sum to
$16+1+16+1+4(4+4)+3(1+1)=72$.
There are no cycles of length one or two, so $\pi(4)=N_4/4=18$. Directly, choose two vertices in each part in $\binom32^2=9$ ways to specify an undirected four-cycle, then choose either orientation. Every walk alternates parts, so no odd-length walk is closed, and all odd $N_k$ and $\pi(k)$ vanish.
\end{solution}

\begin{solution}{ex:v8c07c}
\idea A fixed positive mass near the endpoint survives after taking large roots.
\writeup The upper bound is $m_{2r}\le R^{2r}$. Fix $0<\delta<R$ and write $c_\delta$ for the positive mass of $(R-\delta,R]$. On that interval $x^{2r}\ge(R-\delta)^{2r}$, so
\[
(R-\delta)c_\delta^{1/(2r)}\le m_{2r}^{1/(2r)}\le R.
\]
The fixed factor tends to one. Taking lower limits and then letting $\delta$ decrease to zero proves convergence to $R$.
For the tree measure $R=2\sqrt{d-1}$ and $m_{2r}=a_{2r}$. Dividing by $d$ after taking roots gives $p_{2r}^{1/(2r)}\to2\sqrt{d-1}/d$. At degree three this is $2\sqrt2/3$.
\end{solution}

\begin{solution}{ex:v8c07d}
\idea Split vectors into those unchanged and those negated by swapping sheets.
\writeup Order the cover by sheets. Its adjacency matrix is
\[
\widetilde A=\begin{pmatrix}A_+&A_-\\A_-&A_+\end{pmatrix},
\qquad A=A_++A_-,\quad A_s=A_+-A_-.
\]
The isometries $x\mapsto(x,x)/\sqrt2$ and $x\mapsto(x,-x)/\sqrt2$ have orthogonal images filling the cover's space. The first intertwines $A$ with $\widetilde A$, and the second intertwines $A_s$ with it. Thus orthonormal eigenbases on these subspaces prove the spectral union, including multiplicities.
For a triangle with edges $12,23$ positive and $31$ negative, the cover runs
$(1,0),(2,0),(3,0),(1,1),(2,1),(3,1),(1,0)$,
with no further edges, so it is $C_6$. The original triangle has spectrum $2,-1,-1$. Its signed matrix has characteristic polynomial $(z+2)(z-1)^2$, giving $-2,1,1$. The union is precisely the six-cycle spectrum $2,1,1,-1,-1,-2$.
\end{solution}

\section*{Sources and further reading}

\begin{itemize}
\item \emph{History.} M.~S.~Pinsker, ``On the complexity of a concentrator'', \emph{7th International Teletraffic Conference}, Stockholm, 1973, 318/1--318/4. G.~A.~Margulis, ``Explicit constructions of concentrators'', \emph{Problemy
Peredachi Informatsii} 9:4 (1973), 71--80. Y.~Ihara, ``On discrete subgroups of the two by two projective linear group over $\mathfrak p$-adic fields'', \emph{Journal of the Mathematical Society of Japan} 18 (1966), 219--235.
T.~Sunada, ``L-functions in geometry and some applications'', \emph{Lecture Notes in Mathematics} 1201 (1986), 266--284. H.~Bass, ``The Ihara--Selberg zeta function of a tree lattice'', \emph{International Journal of Mathematics} 3
(1992), 717--797. B.~D.~McKay, ``The expected eigenvalue distribution of a large regular graph'', \emph{Linear Algebra and its Applications} 40 (1981), 203--216. A.~Lubotzky, R.~Phillips and P.~Sarnak, ``Ramanujan graphs'',
\emph{Combinatorica} 8 (1988), 261--277.
\item \emph{Randomness.} M.~Ajtai, J.~Koml\'os and E.~Szemer\'edi, ``Deterministic simulation in LOGSPACE'', \emph{Proceedings of the 19th ACM Symposium on Theory of Computing} (1987), 132--140. D.~Gillman, ``A Chernoff bound for
random walks on expander graphs'', \emph{SIAM Journal on Computing} 27 (1998), 1203--1220.
\item \emph{Recent results.} A.~Marcus, D.~A.~Spielman and N.~Srivastava, ``Interlacing families I: bipartite Ramanujan graphs of all degrees'', \emph{Annals of Mathematics} 182 (2015), 307--325. J.~Friedman, \emph{A Proof of Alon's
Second Eigenvalue Conjecture and Related Problems}, Memoirs of the AMS 195, 2008. J.~Huang, T.~McKenzie and H.-T.~Yau, ``Ramanujan property and edge universality of random regular graphs'', arXiv:2412.20263, 2024 (version 2, 2025), Theorem~1.2 and Corollary~1.3.
J.~Bourgain and A.~Gamburd, ``Uniform expansion bounds for Cayley graphs of $\mathrm{SL}_2(\mathbb F_p)$'', \emph{Annals of Mathematics} 167 (2008), 625--642. The bet: L.~Sloman, ``New proof settles decades-old bet
about connected networks'', \emph{Quanta Magazine}, 18 April 2025.
\item S.~Hoory, N.~Linial and A.~Wigderson, ``Expander graphs and their applications'', \emph{Bulletin of the AMS} 43 (2006), 439--561. This survey develops mixing estimates and the zig-zag construction.
\item A.~Terras, \emph{Zeta Functions of Graphs: A Stroll through the Garden}, Cambridge University Press, 2011. Terras develops the graph zeta function, determinant identities and primitive-cycle counting in detail.
\item G.~Davidoff, P.~Sarnak and A.~Valette, \emph{Elementary Number Theory, Group Theory and Ramanujan Graphs}, Cambridge University Press, 2003. This book gives complete proofs of the number-theoretic LPS construction.
\end{itemize}
